\chapter{Manual de instalación}
\label{chap:manual-instalacion}

\lettrine{E}{ste} apéndice recolle o procedemento completo para poñer en marcha a implementación de referencia
da consola \acrshort{diy} descrita neste traballo. Ofrécense dous camiños, en función do modelo de Raspberry Pi
que se empregue: unha imaxe de NixOS lista para gravar, recomendada e válida só para a Raspberry Pi Zero 2~W,
e unha instalación manual sobre a imaxe de Debian do fabricante, que é a única alternativa para a Raspberry Pi Zero orixinal.


\section{Instalación con NixOS}

Esta é a vía recomendada, pero soamente dispoñible para a Raspberry Pi Zero 2~W. Todo o sistema operativo, incluído o soporte da pantalla, o audio e o propio
emulador, vén xa configurado na imaxe, polo que non hai ningún paso de configuración posterior no dispositivo.

\subsection{Obtención da imaxe}

A imaxe pode descargarse xa construída dende a páxina de \textit{releases} do repositorio, co nome
\texttt{gbeed02.img.zst}, ou construírse na propia máquina empregando a caché binaria sempre que estea dispoñible.
Para construíla só se precisa ter Nix instalado, coa funcionalidade de \textit{flakes} habilitada:

\begin{lstlisting}[language=bash]
nix build github:daniqss/gbeed#installerImages.gbeed02
\end{lstlisting}

Como a imaxe é para a arquitectura \texttt{aarch64-linux}, a construción require unha máquina desa arquitectura
ou, no seu defecto, emulación. Nun sistema NixOS habilítase engadindo á configuración:

\begin{lstlisting}
boot.binfmt.emulatedSystems = ["aarch64-linux"];
\end{lstlisting}

\subsection{Gravado na tarxeta}

Para gravala nunha tarxeta microSD recoméndase descomprimila e gravala empregando \texttt{dd}
seleccionando o dispositivo da tarxeta como destino:

\begin{lstlisting}[language=bash]
zstd -dc gbeed02.img.zst | sudo dd of=/dev/mmcblk0 bs=4M status=progress conv=fsync
\end{lstlisting}

O sistema arrancará o emulador directamente unha vez inserida a tarxeta, pero para poder xogar
é necesario copiar as \acrshortpl{rom} ao dispositivo. Isto pódese facer montando a tarxeta e copiando os ficheiros á ruta \texttt{/home/gbeed/roms}.



\section{Instalación manual sobre Debian}

NixOS non conta con soporte para a Raspberry Pi Zero, polo que a instalación parte da imaxe de
Raspberry Pi OS Bookworm Lite que distribúe o fabricante da pantalla, sobre a que se instalan manualmente
o binario e a configuración descritos a continuación.

\subsection{Cadea de arranque}

Tras o inicio de sesión automático en \texttt{tty1} lévantase un
servidor X no que se executa o emulador, mentres \textit{fbcp} copia continuamente o contido do
\textit{framebuffer} principal (\texttt{fb0}) ao da pantalla \acrshort{spi} (\texttt{fb1}).
Na táboa~\ref{tab:arranque-debian} recóllense as pezas que interveñen en cada paso.

\begin{table}[htbp]
  \centering
  \small
  \begin{tabular}{clp{6.2cm}}
    \hline
    \textbf{Paso} & \textbf{Compoñente} & \textbf{Función} \\
    \hline
    1 & \texttt{config.txt}      & Carga os \textit{overlays} da pantalla \acrshort{spi} e do audio por \acrshort{pwm} \\
    2 & \texttt{getty@tty1}      & Inicio de sesión automático sen contrasinal \\
    3 & \texttt{\textasciitilde/.bash\_profile} & Só en \texttt{tty1}: lanza \textit{fbcp} e o servidor X \\
    4 & \texttt{\textasciitilde/.xinitrc}       & Executa o emulador dentro de X \\
    5 & \textit{fbcp}            & Copia \texttt{fb0} a \texttt{fb1} \\
    6 & \texttt{/etc/asound.conf} & Dirixe a saída de audio ao altofalante \\
    \hline
  \end{tabular}
  \caption{Cadea de arranque da instalación sobre Debian.}
  \label{tab:arranque-debian}
\end{table}

\subsection{Instalación do binario}

O binario para \texttt{armv6l} obtense por compilación cruzada nun contedor, executando na máquina
de desenvolvemento:

\begin{lstlisting}[language=bash]
just cross debian
\end{lstlisting}

O resultado, \texttt{gbeed-debian}, está enlazado contra glibc e co \textit{backend} X11. Unha vez
copiado ao dispositivo, instálase nun directorio do \texttt{PATH}:

\begin{lstlisting}[language=bash]
sudo install -m 755 gbeed-debian /usr/local/bin/gbeed
\end{lstlisting}

\subsection{Configuración do sistema}

Precísanse dous \textit{overlays} de \textit{device tree}, que se descargan da documentación do fabricante
e se copian a \texttt{/boot/firmware/overlays/}:

\begin{itemize}
    \item \texttt{waveshare13}, que crea o dispositivo \texttt{/dev/fb1} correspondente á pantalla.
    \item \texttt{audremap
    18}, que reencamiña a saída de audio ao \acrshort{gpio}~18.
\end{itemize}

Actívanse dende \texttt{/boot/firmware/config.txt}, onde tamén se forza unha saída HDMI virtual de
$480\times480$ píxeles, a resolución da pantalla, aínda que non haxa ningún monitor conectado:

\begin{lstlisting}
dtparam=audio=on
dtparam=spi=on
dtoverlay=waveshare13
dtoverlay=audremap18

# rompe fbcp/DispmanX
# dtoverlay=vc4-kms-v3d

disable_fw_kms_setup=1
hdmi_force_hotplug=1
hdmi_group=2
hdmi_mode=87
hdmi_cvt 480 480 60 6 0 0 0
hdmi_drive=2
\end{lstlisting}


O inicio de sesión automático en \texttt{tty1} configúrase cun \textit{drop-in} de systemd. No directorio
\texttt{/etc/systemd/system/getty@tty1.service.d/} colócase un ficheiro no que a primeira liña
baleira descarta a orde orixinal da unidade antes de substituíla:

\begin{lstlisting}
[Service]
ExecStart=
ExecStart=-/sbin/agetty --autologin <usuario> --noclear %I $TERM
\end{lstlisting}

Xa coa sesión iniciada, \texttt{\textasciitilde/.bash\_profile} lanza o resto da pila gráfica.
A comprobación sobre \texttt{tty1} evita que se intente arrancar X nas sesións abertas por \textit{ssh}:

\begin{lstlisting}[language=bash]
if [ "$(tty)" = "/dev/tty1" ]; then
    export FRAMEBUFFER=/dev/fb0
    fbcp &
    startx
fi
\end{lstlisting}

Finalmente, \texttt{\textasciitilde/.xinitrc} define o que se executa dentro de X. As ordes \texttt{xset}
desactivan o apagado automático da pantalla e \texttt{exec} fai que o servidor X remate cando o faga o
emulador:

\begin{lstlisting}[language=bash]
#!/bin/sh
xset -dpms
xset s off
xset s noblank
exec /usr/local/bin/gbeed
\end{lstlisting}

\subsection{Audio}

O sistema emprega ALSA directamente, sen ningún servidor de son intermedio. Das dúas tarxetas que amosa
\texttt{aplay -l}, a \texttt{0} é a saída HDMI, que non se usa, e a \texttt{1} é a analóxica que chega ao
altofalante grazas ao \textit{overlay} \texttt{audremap18}. Como o emulador abre o dispositivo
predeterminado de ALSA, hai que apuntalo a esa segunda tarxeta creando o ficheiro
\texttt{/etc/asound.conf}:

\begin{lstlisting}
pcm.!default {
    type plug
    slave.pcm "hw:1,0"
}
ctl.!default {
    type hw
    card 1
}
\end{lstlisting}

\subsection{Verificación}

Ante calquera fallo, a pila pódese verificar capa por capa coas ordes da táboa~\ref{tab:verificacion-debian},
o que permite illar en que punto da cadea de arranque está o problema. Convén ter en conta que
\textit{fbcp} depende da librería \texttt{libbcm\_host.so.0}, do paquete \texttt{libraspberrypi0}. Se unha
operación de \texttt{apt} a retira, a pantalla queda en negro aínda que o emulador funcione correctamente.

\begin{table}[htbp]
  \centering
  \footnotesize
  \begin{tabular}{lp{5.2cm}p{3cm}}
    \hline
    \textbf{Capa} & \textbf{Orde} & \textbf{Resultado esperado} \\
    \hline
    Pantalla \acrshort{spi} & \texttt{cat /dev/urandom | sudo tee /dev/fb1} & ruído de cores \\
    \textit{fbcp}           & \texttt{pgrep -a fbcp}                        & devolve un PID \\
    \texttt{libbcm\_host}   & \texttt{ldconfig -p | grep bcm\_host}         & librería listada \\
    Servidor X              & \texttt{pgrep -a Xorg}                        & Xorg en \texttt{:0 vt1} \\
    Audio                   & \texttt{speaker-test -twav -l1}               & escóitase o son \\
    \hline
  \end{tabular}
  \caption{Verificación por capas da instalación sobre Debian.}
  \label{tab:verificacion-debian}
\end{table}
